\section{Scaling Laws 的局限性}

\subsection{训练损失 vs. 下游任务}

\begin{figure}[H]
\centering
\includegraphics[width=0.95\textwidth]{slides-images/slide-029.jpg}
\caption{困惑度 Scaling 很好，但下游任务 Scaling 差异巨大\protect\footnotemark}
\end{figure}
\footnotetext{Slides 第30页。}

Scaling Laws 在\textbf{训练损失}（交叉熵/困惑度）上表现极好，但在\textbf{下游任务}上的预测能力差得多。

\begin{warningbox}{困惑度 $\neq$ 下游性能}
左图显示：不同架构和超参数的模型，其参数量与负对数困惑度几乎完美线性相关——意味着``只要计算量够，什么配置都差不多''。

但右图显示了完全不同的情况：在下游任务（如 SuperGLUE）上，不同架构之间存在\textbf{显著差异}，某些模型远好于其他模型。

\textbf{不要将困惑度的 Scaling 等同于下游能力的 Scaling！}
\end{warningbox}

这一现象在 State Space Models 上也有体现：SSM 在困惑度上可以与 Transformer 匹配甚至超越，但在 in-context learning 和 QA 等能力上可能明显落后。

\subsection{反缩放与分布外行为}

还有一类被称为``反缩放''（inverse scaling）的现象：模型变大后某些行为反而变差。例如，更大的模型更擅长复制训练数据中的模式，因此如果你希望抑制复制行为，这会变得更困难。

\begin{knowledgebox}{反缩放的本质}
反缩放通常出现在\textbf{分布外}（out-of-distribution）区域：当期望的行为与训练数据不一致时，更大的模型更倾向于遵循训练数据的模式，而非人类期望的行为。从某种意义上说，这是经典深度学习鲁棒性问题在大模型时代的延伸。
\end{knowledgebox}

\subsection{本章小结}

Scaling Laws 最可靠的预测对象是训练损失（如交叉熵、困惑度），而非下游任务性能。在使用 Scaling Laws 做工程决策时，需要清楚地区分这两者。对于需要特定能力的应用场景，单独验证是必不可少的。

%% ========== Part 7: 联合 Scaling ==========

